% --- 题目 ---
\problem[P131-8 (16-1)]{已知函数 $\displaystyle f(x)$ 可导, 且 $\displaystyle f(0)=1, 0 < f'(x) < \frac{1}{2}$.
设数列 $\displaystyle \{x_n\}$ 满足 $\displaystyle x_{n+1} = f(x_n) \ (n=1, 2, \dots)$. 证明:
\begin{enumerate}
\item[(1)] 级数 $\displaystyle \sum_{n=1}^{\infty} (x_{n+1} - x_n)$ 绝对收敛;
\item[(2)] $\displaystyle \lim_{n \to \infty} x_n$ 存在, 且 $\displaystyle 0 < \lim_{n \to \infty} x_n < 2$.
\end{enumerate}}
\ansat{316；【十年真题】 - 考点：常数项级数的收敛性 - 8}

% --- 题目 ---
\problem[P131-5 (19-3)]{若 $\displaystyle \sum_{n=1}^{\infty} n u_n$ 绝对收敛, $\displaystyle \sum_{n=1}^{\infty} \frac{v_n}{n}$ 条件收敛, 则 ( \quad )
\begin{tasks}(2)
  \task $\displaystyle \sum_{n=1}^{\infty} u_n v_n$ 条件收敛.
  \task $\displaystyle \sum_{n=1}^{\infty} u_n v_n$ 绝对收敛.
  \task $\displaystyle \sum_{n=1}^{\infty} \left( u_n + v_n \right)$ 收敛.
  \task $\displaystyle \sum_{n=1}^{\infty} \left( u_n + v_n \right)$ 发散.
\end{tasks}}
\ansat{315；【十年真题】 - 考点：常数项级数的收敛性 - 5}

% --- 题目 ---
\problem[P134-例3 (97-1)]{设 $\displaystyle a_1 = 2, a_{n+1} = \frac{1}{2} \left(a_n + \frac{1}{a_n}\right)$ \ ($\displaystyle n = 1, 2, \dots$), 证明:
\begin{enumerate}
\item[(1)] $\displaystyle \lim_{n \to \infty} a_n$ 存在;
\item[(2)] 级数 $\displaystyle \sum_{n=1}^{\infty} \left(\frac{a_n}{a_{n+1}} - 1\right)$ 收敛.
\end{enumerate}}
\ansat{134；【方法探究】 - 考点：常数项级数的收敛性 - 例3}

% --- 题目 ---
\problem[P131-4 (19-1)]{设 $\displaystyle \{u_n\}$ 是单调增加的有界数列, 则下列级数中收敛的是 ( \quad )
\begin{tasks}(4)
  \task $\displaystyle \sum_{n=1}^{\infty} \frac{u_n}{n}.$
  \task $\displaystyle \sum_{n=1}^{\infty} (-1)^n \frac{1}{u_n}.$
  \task $\displaystyle \sum_{n=1}^{\infty} \left(1 - \frac{u_n}{u_{n+1}}\right).$
  \task $\displaystyle \sum_{n=1}^{\infty} \left( u_{n+1}^2 - u_n^2 \right).$
\end{tasks}}
\ansat{315；【十年真题】 - 考点：常数项级数的收敛性 - 4}

% --- 题目 ---
\problem[P134-变式2.2 (94-1,3)]{设常数 $\displaystyle \lambda > 0$, 且级数 $\displaystyle \sum_{n=1}^{\infty} a_n^2$ 收敛, 则级数 $\displaystyle \sum_{n=1}^{\infty} (-1)^n \frac{|a_n|}{\sqrt{n^2 + \lambda}}$ ( \quad )
\begin{tasks}(4)
  \task 发散.
  \task 条件收敛.
  \task 绝对收敛.
  \task 敛散性与 $\displaystyle \lambda$ 有关.
\end{tasks}}
\ansat{316；【方法探究】 - 考点：常数项级数的收敛性 - 变式2.2}

% --- 题目 ---
\problem[P134-变式3 (04-1)]{设有方程 $\displaystyle x^n + nx - 1 = 0$, 其中 $\displaystyle n$ 为正整数. 证明此方程存在唯一正实根 $\displaystyle x_n$, 并证明当 $\displaystyle \alpha > 1$ 时, 级数 $\displaystyle \sum_{n=1}^{\infty} x_n^\alpha$ 收敛.}
\ansat{316；【方法探究】 - 考点：常数项级数的收敛性 - 变式3}

% --- 题目 ---
\problem[P136-6 (21-1)]{设 $\displaystyle u_n(x) = \mathrm{e}^{-nx} + \frac{x^{n+1}}{n(n+1)} \ (n = 1, 2, \cdots)$, 求级数 $\displaystyle \sum_{n=1}^{\infty} u_n(x)$ 的收敛域及和函数.}
\ansat{318；【十年真题】 - 考点二：幂级数的和函数 - 6}

% --- 题目 ---
\problem[P137-9 (17-3)]{设 $\displaystyle a_0 = 1, \ a_1 = 0, \ a_{n+1} = \frac{1}{n+1}(na_n + a_{n-1}) \ (n = 1, 2, 3, \ldots), \ S(x)$ 为幂级数 $\displaystyle \sum_{n=0}^{\infty} a_n x^n$ 的和函数.
\begin{enumerate}
\item[(1)] 证明幂级数 $\displaystyle \sum_{n=0}^{\infty} a_n x^n$ 的收敛半径不小于 $1$;
\item[(2)] 证明 $\displaystyle (1-x)S'(x) - xS(x) = 0 \ (x \in (-1, 1))$, 并求 $\displaystyle S(x)$ 的表达式.
\end{enumerate}}
\ansat{319；【十年真题】 - 考点二：幂级数的和函数 - 9}

% --- 题目 ---
\problem[P137-10 (16-3)]{求幂级数 $\displaystyle \sum_{n=0}^{\infty} \frac{x^{2n+2}}{(n+1)(2n+1)}$ 的收敛域及和函数.}
\ansat{319；【十年真题】 - 考点二：幂级数的和函数 - 10}

% --- 题目 ---
\problem[P137-8 (20-1)]{设数列 $\displaystyle \{a_n\}$ 满足
\[ a_1 = 1, \quad (n+1)a_{n+1} = \left(n+\frac{1}{2}\right)a_n. \]
证明: 当 $\displaystyle |x| < 1$ 时, 幂级数 $\displaystyle \sum_{n=1}^{\infty} a_n x^n$ 收敛, 并求其和函数.}
\ansat{319；【十年真题】 - 考点二：幂级数的和函数 - 8}

% --- 题目 ---
\problem[P136-2 (23-3)]{$\displaystyle \sum_{n=0}^{\infty} \frac{x^{2n}}{(2n)!} = \underline{\hspace{4em}}.$}
\ansat{318；【十年真题】 - 考点二：幂级数的和函数 - 2}

% --- 题目 ---
\problem[P136-4 (17-1)]{幂级数 $\displaystyle \sum_{n=1}^{\infty} (-1)^{n-1} n x^{n-1}$ 在区间 $(-1, 1)$ 内的和函数 $\displaystyle S(x) = \underline{\hspace{4em}}.$}
\ansat{318；【十年真题】 - 考点二：幂级数的和函数 - 4}

% --- 题目 ---
\problem[P140-例(2) (03-1)]{将函数 $\displaystyle f(x) = \arctan \frac{1-2x}{1+2x}$ 展开成 $x$ 的幂级数, 并求级数 $\displaystyle \sum_{n=0}^{\infty} \frac{(-1)^n}{2n+1}$ 的和.}
\ansat{140；【方法探究】 - 考点一：幂级数的收敛域 - 例(2)}

% --- 题目 ---
\problem[P136-2 (20-3)]{设幂级数 $\displaystyle \sum_{n=1}^{\infty} na_n (x-2)^n$ 的收敛区间为 $(-2, 6)$, 则 $\displaystyle \sum_{n=1}^{\infty} a_n (x+1)^{2n}$ 的收敛区间为 ( \quad )
\begin{tasks}(4)
\task $(-2, 6).$
\task $(-3, 1).$
\task $(-5, 3).$
\task $(-17, 15).$
\end{tasks}}
\ansat{317；【十年真题】 - 考点一：幂级数的收敛域 - 2}

% --- 题目 ---
\problem[P136-1 (20-1)]{设 $\displaystyle R$ 为幂级数 $\displaystyle \sum_{n=1}^{\infty} a_n x^n$ 的收敛半径, $\displaystyle r$ 是实数, 则 ( \quad )}
\begin{tasks}(2)
\task 当 $\displaystyle \sum_{n=1}^{\infty} a_{2n} r^{2n}$ 发散时, $\displaystyle |r| \ge R$.
\task 当 $\displaystyle \sum_{n=1}^{\infty} a_{2n} r^{2n}$ 收敛时, $\displaystyle |r| \le R$.
\task 当 $\displaystyle |r| \ge R$ 时, $\displaystyle \sum_{n=1}^{\infty} a_{2n} r^{2n}$ 发散.
\task 当 $\displaystyle |r| \le R$ 时, $\displaystyle \sum_{n=1}^{\infty} a_{2n} r^{2n}$ 收敛.
\end{tasks}
\ansat{317；【十年真题】 - 考点一：幂级数的收敛域 - 1}